\section{An actual large-prime obstruction to central radical partitions}
\label{lp18-section}

Let \(e,f\ge2\) be integers, and set
\[
 c=2^e3^f,\qquad M=2^e,\quad N=3^f,\qquad
 A=M/2,\quad B=N/3,\quad Q=\operatorname{rad}(c-1).
\]
Throughout this section impose the complete balanced two-source hypotheses
\begin{equation}\label{lp-domain}
 \max(M,N)<2\min(M,N)-1,\qquad AB>Q.
\end{equation}
The endpoint \((1,c-1,c)\) is primitive, and its exact source prime
powers are \(M,N\). The no-proper-face theorem and optimization proved
in Section~\ref{dp-section} give
\begin{align}
 \mathcal B_{\rm FCRT}=\mathcal B_{\rm SCRT}=\mathcal B_{\rm PBT}
     &=D+\Gamma,\label{lp-budget}\\
 D&=\log(AB/Q)>0,\notag\\
 \Gamma&=\min_{H\mid Q}
    \operatorname{dist}\bigl(\log H,[\log(Q/B),\log A]\bigr).\notag
\end{align}
Every divisor in this formula is a divisor of the full actual squarefree
integer \(Q\). When a prime is assigned to a packet, its full prime power
on the original endpoint is assigned. No unknown cofactor is omitted.

\begin{theorem}[The exact large-prime partition formula]
\label{lp-exact-barrier}
Under \eqref{lp-domain}, suppose an actual prime \(q\mid c-1\)
satisfies \(q>\max(A,B)\). Put \(R=Q/q\). Then
\[
 \Gamma=\log\frac{q}{\max(A,B)}>0,\qquad
 D+\Gamma=\log\frac{\min(A,B)}R.
\]
No primality condition on \(R\), prescription of the remaining prime
depths, or square condition on \(c\) is required.
\end{theorem}
\begin{proof}
Since \(AB>qR\) and \(q>\max(A,B)\),
\[
 R<\min(A,B),\qquad R<Q/B,\qquad A<q.
\]
Each divisor \(H\mid Q\) either omits \(q\), hence divides \(R\)
and satisfies \(H\le R\), or contains \(q\), hence satisfies
\(H\ge q\). Both extremal choices \(R,q\) are actual divisors.
The central interval is nonempty because \(AB>Q\), and lies strictly
between their logarithms. The distances from its endpoints to the
nearest divisor below and above are therefore
\[
 \log\frac{Q/B}{R}=\log\frac qB,\qquad \log\frac qA.
\]
Their minimum proves the formula for \(\Gamma\). Adding
\(\log(AB/(qR))\) proves the budget formula. This optimizes over all
actual divisors. There is at most one such large prime, since two would
give \(Q>\max(A,B)^2\ge AB\).
\end{proof}

\begin{theorem}[A fully certified nonsquare endpoint]
\label{lp-actual-example}
Take \(e=64,f=41\), and put
\[
 q=3981112602195296746201614890054671463.
\]
Then \(q\) is prime, \eqref{lp-domain} holds, and
\begin{align*}
 c&=672808029771005150108072916419239477248=13^2q+1,\\
 M&=18446744073709551616,
       &N&=36472996377170786403,\\
 A&=9223372036854775808,
       &B&=12157665459056928801,\\
 Q&=51754463828538857700620993570710729019=13q.
\end{align*}
The exact optimum is
\[
 D=\log\frac{c}{78q},\qquad
 \Gamma=\log\frac q{3^{40}},\qquad
 \mathcal B_{\rm FCRT}=\log\frac{2^{63}}{13}.
\]
Moreover \(0<D<1\) and \(\Gamma>52D\).
\end{theorem}
\begin{proof}
Primality is established by the complete recursive certificate proved
below. Direct integer arithmetic gives
\[
 2\min(M,N)-1-\max(M,N)=420491770248316828>0.
\]
The full factorization is \(c-1=13^2q\); thus the four divisors of its
radical are exactly \(1,13,q,13q\). Also
\[
 AB/Q=c/(78q)>1,\qquad q>B>A>13.
\]
The preceding theorem gives all three logarithmic formulas. The exact
integer comparison
\begin{equation}\label{lp-integer-gap-comparison}
 q(78q)^{52}>3^{40}c^{52}
\end{equation}
is checked by the integer certificate, and implies \(\Gamma>52D\)
by strict monotonicity of the logarithm. Finally
\[
 1<\frac{169q+1}{78q}<\frac73<\exp(1)
\]
proves \(0<D<1\). The last inequality follows, for example, from
\(\exp(1)>1+1+1/2>7/3\). No decimal approximation proves any of these
strict inequalities. Since \(f=41\) is odd, \(c\) is not a square,
so this example does not belong to the DP square family.
\end{proof}

For scale only, \(D\) is approximately \(0.77319\), \(\Gamma\)
approximately \(40.33013\), and their sum approximately \(41.10332\).
The exact formulas and comparison \eqref{lp-integer-gap-comparison}
are the mathematical statements.

\begin{corollary}[A scalar-control strengthening fails]
\label{lp-scalar-control-fails}
The assertion that every actual endpoint satisfying
\eqref{lp-domain} has \(\Gamma\le D\) is false, even if one adds
\(0<D<1\). In particular, universal vanishing of \(\Gamma\) on that
class is false.
\end{corollary}
\begin{proof}
The fully factored endpoint of Theorem~\ref{lp-actual-example}
satisfies every premise and \(\Gamma>52D>D>0\).
\end{proof}

This is a finite counterexample to the two displayed strengthenings.
It does not disprove a bound with an unspecified additive constant
\(C_\varepsilon\), the parent FCRT uniform gate, or ABC. All reusable
proper flags are absent by the same inherited theorem used in
\eqref{lp-budget}; no flag or radical factor has been suppressed to
manufacture the failure.

\begin{proposition}[A conditional distinction between two uniform gates]
\label{lp-conditional-infinite}
Fix a prime \(p\ge7\). If there are infinitely many distinct balanced
endpoints \(c=2^e3^f\), \(e,f\ge2\), such that
\[
 c-1=p^2q,\qquad q\text{ prime},
\]
then, along these endpoints as \(c\to\infty\),
\begin{align*}
 D&=\log(p/6)+o(1),\\
 \mathcal B_{\rm FCRT}&=\tfrac12\log c+O_p(1),\\
 \log\operatorname{rad}(c(c-1))&=\log c+O_p(1).
\end{align*}
Consequently the particular budget in \eqref{lp-budget} would fail
its \(\varepsilon\)-log-radical upper bound, with any additive constant,
for each fixed \(0<\varepsilon<1/2\). Meanwhile the scalar ABC excess
\[
 \log c-(1+\varepsilon)\log\operatorname{rad}(c(c-1))
\]
would tend to minus infinity for every fixed \(\varepsilon>0\).
\end{proposition}
\begin{proof}
The distinct positive integers \(c\) tend to infinity along any increasing
enumeration. The balance condition gives \(1/2<M/N<2\); hence \(A,B\)
are comparable to \(\sqrt c\), with absolute constants. Eventually
\(q=(c-1)/p^2\) is distinct from \(p\) and larger than both \(A,B\).
Then \(Q=pq\),
\[
 D=\log\frac{c}{6pq}\longrightarrow\log(p/6)>0.
\]
All hypotheses of Theorem~\ref{lp-exact-barrier} eventually hold, and
it gives
\[
 \mathcal B_{\rm FCRT}=\log\frac{\min(A,B)}p
       =\tfrac12\log c+O_p(1).
\]
Since \(\operatorname{rad}(c(c-1))=6pq=6(c-1)/p\), the radical
asymptotic follows. Subtracting the corresponding multiples of
\(\log c\) proves the two contrasting conclusions.
\end{proof}

No infinitude of these prime values has been proved or used. This
conditional proposition identifies an additional arithmetic input
which would separate this specific certificate gate from scalar ABC.
It must not be reported as an actual infinite counterexample. An
equivalence proved for a different canonical ABC gate does not give
the reverse implication for the present FCRT optimization.

\subsection{A deterministic recursive primality and endpoint certificate}

\begin{lemma}[Full predecessor-factor certificate]
\label{lp-lucas-criterion}
Let \(m>2\) be an integer whose complete factorization of \(m-1\)
into proved primes is known. Suppose that for every prime
\(r\mid m-1\) there is an integer \(a_r\) satisfying
\[
 a_r^{m-1}\equiv1\pmod m,\qquad
 \gcd\bigl(a_r^{(m-1)/r}-1,m\bigr)=1.
\]
Then \(m\) is prime. The witnesses for different primes need not agree.
\end{lemma}
\begin{proof}
Let \(\ell\) be any prime divisor of \(m\). The order of \(a_r\)
modulo \(\ell\) divides \(m-1\), but does not divide \((m-1)/r\).
Its \(r\)-adic valuation is therefore the entire \(r\)-adic valuation
of \(m-1\). This order also divides \(\ell-1\). Applying this to
each prime \(r\mid m-1\) shows \(m-1\mid\ell-1\), whence
\(\ell\ge m\). Since \(\ell\mid m\), we obtain \(m=\ell\).
Induction on \(m\), beginning with the elementary prime two,
justifies a recursive certificate using only smaller prime nodes.
\end{proof}

For the prime in Theorem~\ref{lp-actual-example}, the top decomposition is
\[
 q-1=2\cdot61\cdot563\cdot921889\cdot2532974063
             \cdot24821433536218331.
\]
The witnesses for these six factors are respectively \(5,2,2,2,2,2\).
The complete data contains the predecessor factorizations and witnesses
for all smaller factors, and a node proving thirteen. The standard-library
program \texttt{replay\_large\_prime\_partition.py} checks all twenty-five
prime nodes and seventy modular/gcd witnesses. It uses no factorization
oracle or probable-prime predicate. Thus generation of candidate data
with a computer algebra system leaves no primality-test assumption in
the verifier's proof chain.

After primality, the program checks the exact endpoint factorization,
balance, positive defect, the complete four-divisor list, the rational
identity
\[
 \exp(D)\exp(\Gamma)=\exp(\mathcal B_{\rm FCRT}),
\]
and \eqref{lp-integer-gap-comparison}, using only integers and rational
numbers. Default mode writes a canonical result; \texttt{--check}
recomputes it and compares bytes without changing it. Both modes
have actually returned PASS. The canonical result SHA256 is
\begin{quote}\ttfamily
3b4de60c0f5fb1fc\allowbreak ba68b7def3b655db\allowbreak
4b406047de5ea764\allowbreak f7b29b66f729fb51.
\end{quote}
This certifies one endpoint and its deterministic primality inputs.
The full-class optimization and the conditional implication retain
their ordinary proofs above. Neither the certificate nor this section
adds a Lean theorem, prime-value infinitude, a fixed-\(\varepsilon\)
ABC counterexample, or a proof or disproof of the parent route.
